固定 $t$, 考虑空间谱方程 $\phi_x=U\phi$, 其中 $U=-\ii k\sg+Q$. 在 $z\in\Sigma\setminus\{0,\pm\ii q_0,z_a^\pm\}$ 上定义 Jost 解及修正 Jost 解
\begin{equation}\label{eq:jost-normalization}
\phi_\pm(x,t,z)e^{-\ii\theta(x,t,z)\sg}\longrightarrow Y(z),\qquad x\to\pm\infty,
\end{equation}
\begin{equation}\label{eq:modified-jost}
\mu_\pm=\phi_\pm e^{-\ii\theta\sg},\qquad \mu_\pm\longrightarrow Y\quad(x\to\pm\infty).
\end{equation}
记 $\Delta Q=Q-Q_\infty$. 以下先在分支点之外构造解析列, 再讨论分支点的连续延拓.

\begin{theorem}[Jost 解与解析列]\label{thm:jost}
设 $q-q_0\in L^1(\mathbb R)\cap L^\infty(\mathbb R)$. 对 $z\in\Sigma\setminus\{0,\pm\ii q_0\}$, Volterra 方程
\begin{equation}\label{eq:volterra}
\mu_\pm(x,z)=Y(z)+\int_{\pm\infty}^{x}
Y(z)e^{-\ii\lambda(z)(x-y)\sg}Y(z)^{-1}
\Delta Q(y)\mu_\pm(y,z)e^{\ii\lambda(z)(x-y)\sg}\,\dd y
\end{equation}
有唯一关于 $x$ 有界的解. 列 $\mu_{-,1},\mu_{+,2}$ 在 $D_+$ 解析, 列 $\mu_{+,1},\mu_{-,2}$ 在 $D_-$ 解析, 并连续延拓到各自闭域中避开 $0,\pm\ii q_0$ 的紧集. 在共振原像之外, $\phi_\pm=\mu_\pm e^{\ii\theta\sg}$ 是满足 \eqref{eq:jost-normalization} 的唯一空间 Jost 解.
\end{theorem}
\begin{proof}
修正空间方程及其背景恒等式为
\[
(\mu_\pm)_x=U\mu_\pm+\ii\lambda\mu_\pm\sg,\qquad U_\infty Y=-\ii\lambda Y\sg.
\]
对 $Y^{-1}\mu_\pm$ 使用变分常数公式, 得到 \eqref{eq:volterra}. 令 $s_1=1,s_2=-1$, 第 $j$ 列的传播核为
\begin{equation}\label{eq:columnkernel}
K_j(r,z)=e^{\ii s_j\lambda r}Y e^{-\ii\lambda r\sg}Y^{-1},\qquad r=x-y.
\end{equation}
以 $e_1,e_2$ 表示标准基向量. 对于负端第一列, 令 $w=Y^{-1}\mu_{-,1}$, 则
\begin{equation}\label{eq:jost-column-volterra}
w(x,z)=e_1+\int_{-\infty}^{x}G(x-y,z)\Delta Q(y)Y(z)w(y,z)\,\dd y,
\end{equation}
其中
\begin{equation}\label{eq:jost-Gkernel}
G(r,z)=\diag(1,e^{2\ii\lambda r})Y^{-1}
=\dfrac{1}{\gamma(z)}
\begin{pmatrix}
1&\dfrac{\ii q_0}{z}\\[4pt]
\dfrac{\ii q_0}{z}e^{2\ii\lambda r}&e^{2\ii\lambda r}
\end{pmatrix}.
\end{equation}
当 $y\le x$, $\operatorname{Im}\lambda\ge0$ 时, $|e^{2\ii\lambda(x-y)}|\le1$, 因而此列对应 $D_+$. 对第二列, 相应对角指数为 $\diag(e^{-2\ii\lambda r},1)$; 它在 $y\ge x$, $\operatorname{Im}\lambda\ge0$ 时有界, 给出 $\mu_{+,2}$ 的解析域. 其余两列取相反的积分方向和谱根虚部.

采用向量范数 $\|v\|_1=|v_1|+|v_2|$ 及其诱导矩阵范数. 由
\[
\|\Delta Q(y)\|_1=|q(y)-q_0|,\qquad
\|Y\|_1=1+\dfrac{q_0}{|z|},\qquad
\|Y^{-1}\|_1=\dfrac{1+\dfrac{q_0}{|z|}}{|\gamma(z)|},
\]
可得 HMB 积分核的具体估计
\begin{equation}\label{eq:jost-kernel-estimate}
\|G(x-y,z)\Delta Q(y)Y\|_1
\le c(z)|q(y)-q_0|,\qquad
c(z)=\dfrac{\left(1+\dfrac{q_0}{|z|}\right)^2}{\left|1+\dfrac{q_0^2}{z^2}\right|}.
\end{equation}
取 $\mathcal K\subset\overline{D_+}\setminus\{0,\pm\ii q_0\}$ 为紧集, 记
\[
c_{\mathcal K}=\max_{z\in\mathcal K}c(z),\qquad
M_{\mathcal K}(x)=c_{\mathcal K}\int_{-\infty}^{x}|q(y)-q_0|\,\dd y.
\]
构造 Neumann 级数
\begin{equation}\label{eq:jost-neumann}
\begin{aligned}
w&=\sum\limits_{n=0}^{\infty}w^{(n)},\qquad w^{(0)}=e_1,\\
w^{(n+1)}&=\int_{-\infty}^{x}G(x-y,z)\Delta Q(y)Yw^{(n)}(y,z)\,\dd y.
\end{aligned}
\end{equation}
若 $\|w^{(n)}(y,z)\|_1\le \dfrac{M_{\mathcal K}(y)^n}{n!}$, 则
\begin{align*}
\|w^{(n+1)}(x,z)\|_1
&\le\dfrac{1}{n!}\int_{-\infty}^{x}c_{\mathcal K}|q(y)-q_0|M_{\mathcal K}(y)^n\,\dd y\\
&=\dfrac{M_{\mathcal K}(x)^{n+1}}{(n+1)!}.
\end{align*}
归纳即得
\begin{equation}\label{eq:jost-factorial}
\|w^{(n)}(x,z)\|_1\le\dfrac{M_{\mathcal K}(x)^n}{n!},\qquad
\|w(x,z)-e_1\|_1\le e^{M_{\mathcal K}(x)}-1.
\end{equation}
由于 $M_{\mathcal K}(x)\le c_{\mathcal K}\|q-q_0\|_{L^1}$, 级数对实 $x$ 及 $z\in\mathcal K$ 绝对一致收敛, 且在 $x\to-\infty$ 时趋于 $e_1$. 每一项的积分核在 $D_+$ 解析, 并受上述可积函数控制; 参数积分定理及紧集一致收敛给出 $w$ 和 $\mu_{-,1}=Yw$ 的解析性. 同一控制函数与支配收敛定理给出闭域连续性.

对两个有界解之差 $d$ 迭代 $n$ 次, 有
\[
\|d(x,z)\|_1\le\|d(\cdot,z)\|_{L^\infty}\dfrac{M_{\mathcal K}(x)^n}{n!}\longrightarrow0,
\]
从而解唯一. 正端取 $\displaystyle c_{\mathcal K}\int_x^{+\infty}|q(y)-q_0|\,\dd y$, 反向有序积分给出相同的阶乘估计; 另外两列亦同. 在 $\Sigma$ 上四列均可取边界值, 构成 \eqref{eq:volterra} 的有界矩阵解. 对积分方程作局部绝对连续意义下的 $x$ 微分, 恢复空间谱方程, 其归一化由 \eqref{eq:jost-factorial} 给出. 任一具有同一振荡归一化的空间解通过变分常数公式满足该 Volterra 方程, 因而也唯一.
\end{proof}

\begin{proposition}[向分支点的连续延拓]\label{prop:jost-branch}
在定理 \ref{thm:jost} 的条件上, 再设
\begin{equation}\label{eq:jost-weighted}
\int_{-\infty}^{+\infty}(1+|x|)|q(x)-q_0|\,\dd x<\infty.
\end{equation}
则每个 Jost 解析列均可从其闭解析域连续延拓到 $z_b=\pm\ii q_0$. 端点值满足
\begin{equation}\label{eq:jost-branch-volterra}
\mu_\pm(x,z_b)=Y(z_b)+\int_{\pm\infty}^{x}
[I+(x-y)U_\infty(z_b)]\Delta Q(y)\mu_\pm(y,z_b)\,\dd y,
\end{equation}
并在对应归一化端趋于 $Y(z_b)$. 此矩阵在任一对应半直线上有界, 满足空间谱方程, 且秩为一.
\end{proposition}
\begin{proof}
由 $U_\infty^2=-\lambda^2I$, 背景传播矩阵可写为
\begin{equation}\label{eq:background-propagator}
e^{rU_\infty}
=\cos(\lambda r)I+\dfrac{\sin(\lambda r)}{\lambda}U_\infty
\longrightarrow I+rU_\infty(z_b),\qquad z\to z_b.
\end{equation}
尽管 $Y^{-1}$ 在分支点发散, \eqref{eq:columnkernel} 中的乘积具有可控的连续极限. 具体地,
\begin{equation}\label{eq:branch-columnkernel}
K_j(r,z)=\dfrac{1+e^{2\ii s_j\lambda r}}{2}I
+\dfrac{e^{2\ii s_j\lambda r}-1}{2\ii s_j\lambda}U_\infty(z).
\end{equation}
在该列允许的积分方向及闭解析域内, $s_j\operatorname{Im}\lambda\,r\ge0$, 且
\[
\dfrac{e^{2\ii s_j\lambda r}-1}{2\ii s_j\lambda}
=r\int_0^1 e^{2\ii s_j\tau\lambda r}\,\dd\tau,\qquad
\|K_j(r,z)\|_1\le1+C_b|r|.
\]
这里 $C_b$ 可在 $z_b$ 的一个闭邻域上一致选取. 这同时给出 $K_j(r,z_b)=I+rU_\infty(z_b)$.

以负端为例, 固定 $b\in\mathbb R$, 对 $y\le x\le b$ 有
\[
1+C_b|x-y|\le A_b(1+|y|),\qquad
L_b(x)=A_b\int_{-\infty}^{x}(1+|y|)|q(y)-q_0|\,\dd y.
\]
直接对 $\mu_{-,j}$ 的列积分方程迭代, 其第 $n$ 项满足
\[
\|\mu_{-,j}^{(n)}(x,z)\|_1
\le\left(\sup_{z}\|Y_j(z)\|_1\right)\dfrac{L_b(x)^n}{n!},\qquad x\le b.
\]
加权可积条件保证级数在 $x\le b$ 及分支点邻域内一致收敛. 各项利用 \eqref{eq:branch-columnkernel} 支配收敛, 故其和连续趋于 \eqref{eq:jost-branch-volterra}; $L_b(x)\to0$ 则给出负端归一化. 正端使用 $y\ge x\ge b$ 和从 $x$ 到 $+\infty$ 的加权积分, 得到相同结论. 对差的有序积分估计还给出端点 Volterra 解在对应半直线上的唯一性.

最后, $U_\infty(z_b)Y(z_b)=0$, $U_\infty(z_b)^2=0$. 对 \eqref{eq:jost-branch-volterra} 求导得到 $(\mu_\pm)_x=U(z_b)\mu_\pm$. 因 $Y(z_b)$ 的两列线性相关, 唯一性使同一关系对 $\mu_\pm(x,z_b)$ 成立. 其归一化端值非零, 而一阶线性系统的非零解不能在有限点为零, 因而矩阵秩恒为一.
\end{proof}

共振原像 $z_a^\pm$ 不使空间积分核退化, 因而 $\mu_\pm$ 在此仍由 \eqref{eq:volterra} 正则取值. 传播相位在共振处的处理及两组 Jost 解的共同演化见第 \ref{ch:scattering} 节.
